% A manuscript fragment. Requires amsmath, amssymb and hyperref.
\section{From primitive characters to the Riemann hypothesis}
\label{sec:character-to-rh}

Let \(N\) be a positive integer. A complex Dirichlet character of level
\(N\) is a character of \((\mathbb Z/N\mathbb Z)^\times\), extended by
zero to residue classes that are not units. For \(\Re s>1\) write
\[
       L(s,\chi)=\sum_{n=1}^{\infty}\frac{\chi(n)}{n^s}.
\]
Elsewhere \(L(s,\chi)\) denotes its analytic continuation at the points
where that continuation is regular. Write \(\mathbf 1_N\) for the
principal character of level \(N\). The all-character assertion used
in the last step of the proof is
\begin{equation}
\label{eq:all-character-nv}
   \Re s>\tfrac12,\qquad
   (\chi\ne\mathbf 1_N\ \text{or}\ s\ne 1)
       \quad\Longrightarrow\quad L(s,\chi)\ne0
\end{equation}
for every positive \(N\) and every character \(\chi\) of level \(N\).
The exception in the principal case removes the pole. Keeping this
disjunction also permits evaluation of a nonprincipal character at
\(s=1\). The primitive nonvanishing assertion in the source is phrased
using a total complex-valued function, including at exceptional points.
Its restriction to the regular points above suffices for the
classical argument.\footnote{\href{https://github.com/Deicyde/riemann-hypothesis/blob/40e56bab/src/Development/AutomorphicArithmetic/CharacterBoxes/EveryCharacter.lean\#L13-L18}{Definition of the all-character assertion and its deduction from the primitive assertion},
\href{https://github.com/Deicyde/riemann-hypothesis/blob/40e56bab/src/Formalism/SharpRestriction/RamificationFrobenius/TotientGaloisClasses/ExplicitFormulaErrorOfHostConditions.lean\#L11-L18}{the primitive assertion}.}

\subsection{Restoring the Euler factors at the level}
\label{sec:restore-euler}

Assume the right-half-plane nonvanishing produced by the primitive-character
argument, and fix \(N,\chi,s\) satisfying the hypotheses in
\eqref{eq:all-character-nv}. Put \(f=\operatorname{cond}(\chi)\).
Then \(f\) divides \(N\), and there is a primitive character \(\chi^*\)
of level \(f\) inducing \(\chi\). If \(s=1\) we have
\(\chi\ne\mathbf 1_N\), hence \(\chi^*\ne\mathbf 1_f\), since raising
the level of a principal character yields the principal character.
Thus the primitive character is evaluated at a regular point in all
cases under consideration. The primitive nonvanishing statement gives
\(L(s,\chi^*)\ne0\).

The change-of-level identity in the normalization used here is
\begin{equation}
\label{eq:dirichlet-change-level}
     L(s,\chi)=L(s,\chi^*)\prod_{p\mid N}
                     \bigl(1-\chi^*(p)\,p^{-s}\bigr),
\end{equation}
where \(p\) ranges over distinct rational primes. For \(p\mid f\) the
factor is \(1\), since \(\chi^*(p)=0\). For \(\Re s>1\) the identity
follows by deleting the Euler factors at primes dividing \(N\) from
the series for \(\chi^*\). Holomorphy gives the identity on
\(\mathbb C\setminus\{1\}\) by analytic continuation. If \(s=1\) is
allowed, both characters are nonprincipal and continuity gives the same
identity there.\footnote{\href{https://github.com/leanprover-community/mathlib4/blob/v4.33.1/Mathlib/NumberTheory/LSeries/DirichletContinuation.lean\#L125-L162}{Change of level and extension to \(s=1\).}}

For every \(p\mid N\), the value of \(\chi^*(p)\) is zero or a root of
unity, so \(\lvert\chi^*(p)\rvert\le1\). Also
\[
       \lvert\chi^*(p)\,p^{-s}\rvert
       \le p^{-\Re s}<1.
\]
The strict inequality follows from \(p>1\) and \(\Re s>\tfrac12\).
Consequently no factor in \eqref{eq:dirichlet-change-level} vanishes.
Multiplication by the nonzero value \(L(s,\chi^*)\) proves
\eqref{eq:all-character-nv}.\footnote{\href{https://github.com/Deicyde/riemann-hypothesis/blob/40e56bab/src/Development/AutomorphicArithmetic/CharacterBoxes/Imprimitive.lean\#L5-L35}{Reduction from arbitrary characters to primitive characters.}}

\subsection{The character of level one}
\label{sec:level-one-specialization}

The group \((\mathbb Z/\mathbb Z)^\times\) is trivial. Its unique
character takes the value \(1\) on every positive integer. Its
\(L\)-function therefore agrees with \(\zeta(s)\) for \(\Re s>1\),
and their continuations agree away from the pole. Specializing
\eqref{eq:all-character-nv} to \(N=1\) and \(\chi=\mathbf 1_1\)
gives
\begin{equation}
\label{eq:zeta-right-half}
             \zeta(s)\ne0
             \qquad(\Re s>\tfrac12,\ s\ne1).
\end{equation}
Both restrictions in \eqref{eq:zeta-right-half} are exactly those used in
the reflection argument.\footnote{\href{https://github.com/Deicyde/riemann-hypothesis/blob/40e56bab/src/Formalism/GammaFactorNonvanishing/CritLine.lean\#L29-L32}{Specialization to the character modulo one.}}

\subsection{Reflection, including the possible zeros of the cosine}
\label{sec:zeta-reflection}

\noindent\textbf{Proposition.}
Suppose \eqref{eq:zeta-right-half} holds. For every \(s\in\mathbb C\),
if \(\zeta(s)=0\), \(s\ne1\), and \(s\ne-2(n+1)\) for each
\(n\in\mathbb Z_{\ge0}\), then \(\Re s=\tfrac12\).

\medskip
\noindent\textit{Proof.}
Fix such a zero \(s\). The possibility \(\Re s>\tfrac12\) is excluded
by \eqref{eq:zeta-right-half}. Suppose then that \(\Re s<\tfrac12\),
and set \(w=1-s\). Since \(\zeta(0)=-\tfrac12\), we have \(s\ne0\),
and hence \(w\ne1\). We also have \(\Re w>\tfrac12\), which puts
\(w\) outside \(\{0,-1,-2,\dots\}\). These statements justify using
the zeta functional equation at \(w\), which gives
\begin{equation}
\label{eq:zeta-cos-reflection}
    \zeta(1-w)=2(2\pi)^{-w}\Gamma(w)
                       \cos\!\left(\frac{\pi w}{2}\right)\zeta(w).
\end{equation}
The complex power \((2\pi)^{-w}\) is nonzero because \(2\pi>0\).
The gamma factor is nonzero because \(\Re w>0\). Equation
\eqref{eq:zeta-right-half} at \(w\) gives \(\zeta(w)\ne0\).

It remains to test the cosine. If it vanished, the complex zeros of
cosine would give \(w=2k+1\) for some \(k\in\mathbb Z\), and
therefore \(s=-2k\). If \(k<0\), then \(s\ge2\), contradicting
\(\Re s<\tfrac12\). If \(k=0\), then \(s=0\), already excluded.
If \(k>0\), take \(n=k-1\) to contradict \(s\ne-2(n+1)\).
Thus every factor on the right of \eqref{eq:zeta-cos-reflection} is
nonzero, whereas the left side is \(\zeta(s)=0\). This excludes
\(\Re s<\tfrac12\) and proves the proposition.
\hfill\(\square\)

This proposition has precisely the quantifiers used for the classical
Riemann hypothesis in the released theorem. Its nonvanishing premise
follows from the all-character result through
\eqref{eq:zeta-right-half}.\footnote{\href{https://github.com/Deicyde/riemann-hypothesis/blob/40e56bab/src/Formalism/GammaFactorNonvanishing/RhOfRight.lean\#L5-L62}{The reflection argument},
\href{https://github.com/leanprover-community/mathlib4/blob/v4.33.1/Mathlib/NumberTheory/LSeries/RiemannZeta.lean\#L175-L186}{the zeta functional equation and the formal statement of RH},
\href{https://github.com/Deicyde/riemann-hypothesis/blob/40e56bab/src/Supplements/Kernels/MainTheorems.lean\#L7-L9}{application of the all-character result},
\href{https://github.com/Deicyde/riemann-hypothesis/blob/40e56bab/src/Main.lean\#L4-L6}{the two classical RH declarations}.}

\section{How the arithmetic input supplies the character theorem}
\label{sec:entry-assembly}

The all-character assertion \eqref{eq:all-character-nv} is obtained by
applying a theorem with two inputs. One input includes the existence of
local Langlands parameters in the normalization described below, together
with two uses of a local--global compatibility statement. The other
input supplies arithmetic windows for the auxiliary fields. The two
uses of local--global compatibility are instantiated by exactly the same
theorem. Although its first input lists three components, the proof
has to establish two different kinds of assertion to fill them.
\footnote{\href{https://github.com/Deicyde/riemann-hypothesis/blob/40e56bab/src/Development/AutomorphicArithmetic/ThetaCorrespondence/ThetaPairLifts/ThirdPlaceSeries/EveryCharacterOfLiteratureOfVWin.lean\#L210-L218}{The three components of the input},
\href{https://github.com/Deicyde/riemann-hypothesis/blob/40e56bab/src/Development/AutomorphicArithmetic/CharacterBoxes/MainConclusion.lean\#L6-L12}{their instantiation in the final argument}.}

\subsection{The local parameter used in the comparison}

Fix a number field \(E\), an algebraic closure \(\Omega\), a finite
place \(w\) of \(E\), and a place of \(\Omega\) above \(w\). Put
\(K=E_w\). For each \(n\ge1\), consider irreducible smooth complex
representations of \(\mathrm{GL}_n(K)\) up to isomorphism. The local
construction supplies a parameter \(r(\pi)\) of Weil--Deligne type,
together with a local Artin map \(a\) from the Weil group to \(K^\times\).
On characters it gives \(r(\chi)=\chi\circ a\) with zero monodromy. It
sends cuspidal representations to irreducible parameters and sends a
Langlands quotient for a multisegment to the direct sum of the parameters
of its segments. It preserves dimension and intertwines real powers of
the norm. For every \(\pi\) with central character \(\omega_\pi\), it
also satisfies
\[
   |\det r(\pi)(\sigma)|=|\omega_\pi(a(\sigma))|
\]
on the Weil group, with \(\det\) taken on the Weil action. Several
stronger assertions in this construction have a restriction on the
cuspidal support. An allowed cuspidal constituent has rank one, or is a
smooth-character twist of a representation with first principal
congruence invariants over a field topologically isomorphic to
\(\mathbb Q_p\). If \(\pi\) has a cuspidal support all of whose members
are allowed, \(r\) respects smooth character twists, the central
character equality holds without absolute values, and duals correspond.
For such a \(\pi\), equality \(r(\pi')=r(\pi)\) implies \(\pi'=\pi\).
If in addition \(\pi\) has first principal congruence invariants, its
parameter is tame.

More explicitly, suppose \(K\simeq\mathbb Q_g\) and choose an
extension place \(w'\) unramified over \(g\) with residue field
\(\mathbb F_{g^n}\). Suppose the restriction of the local representation
to the integral subgroup contains the inflation of a cuspidal type
of \(\mathrm{GL}_n(\mathbb F_g)\) with regular parameter
\(\theta:\mathbb F_{g^n}^{\times}\to\mathbb C^\times\).
The character of this type at a matrix with irreducible characteristic
polynomial and an eigenvalue \(z\in\mathbb F_{g^n}^{\times}\) is
\((-1)^{n-1}\sum_{j=0}^{n-1}\theta(z^{g^j})\). The monodromy of
\(r(\pi)\) is zero. If \(t\) belongs to inertia at \(w'\) and
\(z\) is the residue of its local Artin image there, regard \(t\)
as an element of the inertia group over \(E\). Its characteristic
polynomial under \(r(\pi)\) is
\[
       \prod_{j=0}^{n-1}\bigl(X-\theta(z^{g^j})\bigr).
\]
\footnote{\href{https://github.com/Deicyde/riemann-hypothesis/blob/40e56bab/src/Development/AutomorphicArithmetic/Foundations/LocalLanglands/IsRecTOfIsRec.lean\#L14-L80}{The restrictions on cuspidal support, the inertia polynomial and the defining properties of the parameter.}}

Here is the extension step used in constructing \(r\).
Choose the Artin map, and first construct parameters \(r_{\rm c}\)
on cuspidal representations. Multisegment exhaustion and the extension
lemma give a map \(r\) on all irreducible representations agreeing with
\(r_{\rm c}\) on cuspidal ones. Suppose \(\pi\) is represented by a
multisegment and let \(S\) be the multiset of its cuspidal terms,
including their norm twists. For \(\sigma\) in the Weil group let
\(P_{D,\sigma}(X)\) denote the characteristic polynomial of the Weil
action on \(D\). The segment calculation gives
\begin{equation}
\label{eq:entry-segment-polynomial}
      P_{r(\pi),\sigma}(X)
      =\prod_{\tau\in S}P_{r_{\rm c}(\tau),\sigma}(X).
\end{equation}
Indeed for a segment of length \(d\) with first term \(\rho\), the Weil
eigenvalues of the segment parameter collect the eigenvalues of the
successive norm twists of \(r_{\rm c}(\rho)\) at exponents
\(0,1,\ldots,d-1\). Norm-twist compatibility on cuspidals identifies
these with the parameters of the terms of the segment. Direct summation
multiplies characteristic polynomials, giving
\eqref{eq:entry-segment-polynomial} for the union of the segments.
At \(\sigma=1\), comparison of degrees in this identity gives
\(\dim r(\pi)=\sum_{\tau\in S}\dim r_{\rm c}(\tau)=n\). For general
\(\sigma\), comparison of products of eigenvalues gives
\(\det r(\pi)(\sigma)=\prod_{\tau\in S}\det r_{\rm c}(\tau)(\sigma)\).
The central character of \(\pi\) is the product of the central
characters of its cuspidal support. Taking absolute values therefore
gives the stated central character identity for every \(\pi\). If
all its constituents are allowed, use the exact central character
identity on each factor to obtain the exact identity for \(\pi\).
\footnote{\href{https://github.com/Deicyde/riemann-hypothesis/blob/40e56bab/src/Development/AutomorphicArithmetic/Foundations/LocalLanglands/LocalCentralUnits/ExtDim.lean\#L19-L78}{The segment identity, its degree and the two determinant conclusions.}}

For a one-dimensional representation, the cuspidal construction gives
the asserted formula from reciprocity directly. Each segment starts at
a cuspidal representation, so the multisegment extension rule with
\(r_{\rm c}\) also gives the required rule with \(r\) in its place.
The extension lemmas for twists, duals, injectivity and depth give
the corresponding properties above. For the inertia assertion the
representation containing the inflated finite cuspidal type is itself
cuspidal. Hence \(r\) agrees with \(r_{\rm c}\) on that representation,
and the inertia calculation for \(r_{\rm c}\) applies. This supplies
the local existence assertion at the chosen place, and the choices of
\(E,\Omega,w\) and a place over \(w\) were arbitrary.\footnote{\href{https://github.com/Deicyde/riemann-hypothesis/blob/40e56bab/src/Development/AutomorphicArithmetic/Foundations/LocalLanglands/LocalCentralUnits/RecTExists.lean\#L20-L66}{Construction of the local parameter and verification of its required laws.}}

\subsection{The global compatibility statement}

Let \(E\) be a CM field with complex conjugation \(c\), and let \(\ell\)
be a rational prime. Fix an algebraic closure \(\Omega\) of \(E\) and
a field isomorphism \(\iota:\overline{\mathbb Q}_{\ell}\to\mathbb C\).
For \(a\ge1\), take a regular algebraic cuspidal representation \(\pi\)
of \(\mathrm{GL}_a\) over \(E\), together with a Hecke character of
its maximal totally real subfield polarizing \(\pi\). The global
statement provides a continuous semisimple representation
\[
  R:\operatorname{Gal}(\Omega/E)\longrightarrow
       \mathrm{GL}_a(\overline{\mathbb Q}_\ell),
\]
an integer \(b\), and a multiset \(H_\tau\) of integers for every
embedding \(\tau:E\to\overline{\mathbb Q}_{\ell}\). If \(I_{\iota\tau}\)
denotes the multiset of infinitesimal-character parameters for \(\pi\)
at \(\iota\tau\), the equalities of multisets are
\[
  H_\tau=\{(a-1)/2-z:z\in I_{\iota\tau}\},
  \qquad H_{\tau c}=\{b-h:h\in H_\tau\}.
\]
The first equality uses the natural inclusion \(\mathbb Z\subset\mathbb C\).

At a finite place \(w\nmid\ell\), let \(x\) be the Weil--Deligne
parameter of \(R\) and \(y\) the local parameter of the algebraically
normalized component of \(\pi\), using the chosen \(\iota\) to compare
coefficients. Their dimensions agree. If all Weil commutators act
trivially in \(y\) and its monodromy \(N_y\) vanishes, then \(x=y\).
There is a second comparison when \(E_w\simeq\mathbb Q_g\) for a
prime \(g\) and either \(a=1\) or \(16a^2<g\). Suppose \(y\) is tame,
\(N_y=0\), and for each inertia element generating tame inertia
topologically modulo wild inertia every repeated eigenvalue \(z\) of
that element in \(y\) satisfies \(z^{\operatorname{N}w-1}=1\). Then
\(N_x=0\) and the characteristic polynomials of \(x\) and \(y\)
agree on every Weil element of degree zero.

At \(w\mid\ell\), with the compatible local period data of the
comparison theorem, a spherical local component gives a crystalline,
de Rham representation with jumps \(H_\tau\) at the embeddings it
induces. If \(N_y=0\) and either the commutator condition or the tame,
repeated-eigenvalue and prime-size conditions above hold, the representation
is de Rham and potentially semistable, with the same prescribed jumps.
Its potentially semistable Weil--Deligne parameter then satisfies the
two comparisons just stated for \(x\) and \(y\). These clauses specify
the local conclusions used in both occurrences of the global input.
\footnote{\href{https://github.com/Deicyde/riemann-hypothesis/blob/40e56bab/src/Development/AutomorphicArithmetic/Foundations/GaloisRepresentations/TypeConditions/WeightsConjBis.lean\#L15-L86}{The global comparison, local hypotheses and weights},
\href{https://github.com/Deicyde/riemann-hypothesis/blob/40e56bab/src/Development/AutomorphicArithmetic/Foundations/GaloisRepresentations/TypeConditions/WeightsConjAlt.lean\#L22-L33}{the commutator and inertia conditions}.}

\subsection{The arithmetic window and its quantifiers}

For a primitive Dirichlet character \(\psi\), \(0<\epsilon<1/52\) and
\(U>1\), the pair of auxiliary primes \(q,d\) satisfies
\begin{align*}
 q&\equiv3\pmod8,& \tfrac12 U^{\epsilon/5}
                   &<\log q\le U^{\epsilon/5},\\
 d&\equiv7\pmod{16},&
 \tfrac12 U^{\epsilon/5}(\log U)^2
                   &<\log d\le U^{\epsilon/5}(\log U)^2 .
\end{align*}
In addition \(q\mid d+1\) and every odd prime \(p\le U^{\epsilon/5}\)
divides \(d+1\). Choose any function \(e:\mathbb R\to\mathbb Z_{\ge0}\)
such that for all \(x\ge1\),
\[
       x^{2+8\epsilon}\le 2^{e(x)}\le 2x^{2+8\epsilon}.
\]
Such a function exists. For \(x\ge1\), take \(n\ge0\) with
\(2^n\le x^{2+8\epsilon}<2^{n+1}\), and put \(e(x)=n+1\).
For \(x<1\) put \(e(x)=0\).

For fixed \(e,q,d\), the number fields \(E\) in the assertion can be
described by a subgroup of a cyclotomic Galois group. There must be
a positive integer \(M\) divisible by \(2^{e+2}\), the conductor of
\(\psi\), \(q\), and \(d\), and a \(\mathbb Q\)-embedding
\(E\hookrightarrow\mathbb Q(\zeta_M)\). Its image is the fixed field
of the automorphisms \(\zeta_M\mapsto\zeta_M^u\) satisfying
\[
  u\equiv\pm1\bmod 2^{e+2},\qquad \psi(u)=1,\qquad
  \Bigl(\frac uq\Bigr)=\Bigl(\frac ud\Bigr)=1.
\]

Let \(S_\ell\) consist of all places of \(E\) dividing \(\ell\),
with a choice of a place of an algebraic closure above each one.
Use the Galois group \(G_{E,S_\ell}\) unramified outside \(S_\ell\)
at finite places, and the discrete module
\(A_{\ell,k}=(\mathbb Q_\ell/\mathbb Z_\ell)(-k)\), on which Galois
acts by the inverse \(k\)-th power of the cyclotomic character.
Let \(D_v\) be the decomposition group at the chosen place over \(v\).
Put
\[
   T_{\ell,k}(E)=\ker\!\left[
      H^1(G_{E,S_\ell},A_{\ell,k})\longrightarrow
           \prod_{v\in S_\ell}H^1(D_v,A_{\ell,k})\right],
   \qquad b_{\ell,k}=\operatorname{length}_{\mathbb Z_\ell} T_{\ell,k}(E).
\]
All cohomology here uses continuous cochains.
\footnote{\href{https://github.com/Deicyde/riemann-hypothesis/blob/40e56bab/src/Formalism/GaloisCohomology/Statement.lean\#L48-L78}{The places, cyclotomic character, sign of the twist and strict cohomology lengths}.}

For this \(\psi,\epsilon,e(\cdot)\), the window assertion supplies
\(J\ge1\) and \(X_0\) such that for every \(X\ge X_0\) and every
pair \(q,d\) satisfying the conditions above at \(U=X^{1/J}\),
one can choose \(Y\ge X^{1/J}\) with \(Y^2\le X\). For \emph{every}
\(Y\le x\le Y^2\), take \(t=\lfloor x^{401}\rfloor\). For every field
\(E\) defined above at \(e=e(x)\), and every choice of algebraic closure
and local data, all \(b_{\ell,k}\) with \(1\le k<t\) are finite.
Simultaneously for every finite set \(P\) of rational primes,
\begin{equation}\label{eq:entry-strict-length}
      \sum_{\ell\in P}(\log\ell)\sum_{k=1}^{t-1}b_{\ell,k}
                 \ \le\ \sqrt{\exp(x^\epsilon)} .
\end{equation}
In particular \(Y\) is chosen before \(x,E,P\). The scale for
choosing the auxiliary primes is \(X^{1/J}\), which may be smaller
than \(Y\). The source formulates the bound over all finite \(P\),
so its use never requires a prior finiteness assertion about the set
of contributing primes.
\footnote{\href{https://github.com/Deicyde/riemann-hypothesis/blob/40e56bab/src/Development/AutomorphicArithmetic/CharacterBoxes/ExistsPageN.lean\#L12-L49}{Definitions of the field, the prime pair and the window inequality},
\href{https://github.com/Deicyde/riemann-hypothesis/blob/40e56bab/src/Development/AutomorphicArithmetic/CharacterBoxes/ExistsPageN.lean\#L67-L77}{existence of the scale function}.}

\subsection{Where the two inputs are used}

On the global side, the geometric realization and the theta-lift
theorem first give the local--global comparison for this normalization
of the parameters. More precisely the geometric input supplies the
initial representations at the odd levels in the construction, while
weight existence supplies the congruences. The passage to the limit,
descent, and the comparison of traces give the simultaneous core
statement. From it one obtains spherical compatibility on the unlinked
support and the Iwahori-fixed comparison. Applying these to the local
and global parameters produces the form of the comparison required
above.\footnote{\href{https://github.com/Deicyde/riemann-hypothesis/blob/40e56bab/src/Development/AutomorphicArithmetic/Foundations/GaloisRepresentations/IntegralStructures/Development.lean\#L18-L31}{Local--global comparison from geometric and theta-lift inputs},
\href{https://github.com/Deicyde/riemann-hypothesis/blob/40e56bab/src/Development/AutomorphicArithmetic/Foundations/GaloisRepresentations/IntegralStructures/RowSrc.lean\#L4-L8}{the resulting Galois theorem}.}

Use this compatibility in the first and third components of the first
input. Between them use the local parameter constructed above. To
obtain the window input, combine the same compatibility with isobaric
base change. The four sets of data used in the window theorem are local
representation theory, automorphic results with base change and Galois
compatibility, Galois cohomology, and the CM and abelian-variety
results. The local data include multisegment classification and the
finite-field cuspidal characters. The CM data include deformation,
comparison of periods, good reduction, the elliptic CM theorem and the
Hecke--Tate functional equation. With the four sets supplied the
window theorem applies. Both inputs to the character theorem have now
been specified on the same Galois comparison.\footnote{\href{https://github.com/Deicyde/riemann-hypothesis/blob/40e56bab/src/Development/AutomorphicArithmetic/AutomorphicVarieties/ModularityArgument/ThresholdArguments/DenseThresholdCases/TopMainZj.lean\#L18-L74}{Local, automorphic, cohomological and CM inputs to the window},
\href{https://github.com/Deicyde/riemann-hypothesis/blob/40e56bab/src/Development/AutomorphicArithmetic/AutomorphicVarieties/ModularityArgument/ThresholdArguments/DenseThresholdCases/VwinOfRow.lean\#L4-L5}{application to the specified compatibility theorem}.}

The final arithmetic implication uses the theta integral equality and
local matching as two separate inputs. It also uses the covolume
estimate for the Hodge bundles, local Langlands existence and the
arithmetic windows just constructed. In the passage from the theta
inputs to the family of comparisons, the proof obtains two lower
thresholds \(C_0\) and \(C_1\). The first ensures the existence of
the required row functionals and the second ensures the evaluation
property on the grid of parameters. Choosing \(C=\max(C_0,C_1)\)
permits using both conclusions for one grid. The linked-bundle
covolume estimate applies with this same \(C\), as does the bound
for the contributions at two and three on the chosen family. These
give the required family of hypotheses. The primitive-character
argument and the Euler-factor calculation in
Section~\ref{sec:restore-euler} now yield
\eqref{eq:all-character-nv}.\footnote{\href{https://github.com/Deicyde/riemann-hypothesis/blob/40e56bab/src/Development/AutomorphicArithmetic/ThetaCorrespondence/ThetaPairLifts/ThirdPlaceSeries/EveryCharacterOfLiteratureOfVWin.lean\#L810-L813}{Inputs to the last arithmetic implication},
\href{https://github.com/Deicyde/riemann-hypothesis/blob/40e56bab/src/Development/AutomorphicArithmetic/ThetaCorrespondence/ThetaPairLifts/ThirdPlaceSeries/SiegelEisensteinSecondTermHypotheses.lean\#L44-L73}{choice of the common threshold and application to characters}.}
